\documentclass[11pt,letterpaper]{article}

\usepackage[margin=0.6in]{geometry}
\usepackage{enumitem}
\usepackage{amsmath}
\usepackage{xcolor}

\pagestyle{empty}
\setlength{\parindent}{0pt}

% Answer block: indented, set in a contrasting color
\newenvironment{answer}{\par\vspace{4pt}\color{blue!60!black}\textbf{Answer:}\ }{\par\vspace{4pt}}

\begin{document}

\begin{center}
    {\Large \textbf{16.07 Dynamics --- Quiz 4: Kane Damper (Solutions)}}\\[4pt]
\end{center}

\vspace{4pt}
The questions below refer to the rigid body with a Kane damper from Problem Set 4. Recall that the fluid applies a torque $C_d\left({}^{B}\boldsymbol{\omega}_d - {}^{B}\boldsymbol{\omega}\right)$ to the body, and that the total kinetic energy changes at the rate
\[
\dot T = -C_d\,\left\lVert {}^{B}\boldsymbol{\omega}_d - {}^{B}\boldsymbol{\omega} \right\rVert^2 .
\]

\vspace{10pt}
\begin{enumerate}[label=\textbf{\arabic*.}, leftmargin=*, itemsep=4pt]

\item Without a damper, spin about which principal axes is stable? Which of these does the Kane damper make unstable, and why?

\begin{answer}
\textbf{Without a damper, both the major and minor axes are stable. The damper makes minor-axis spin unstable.} The damper torques are internal, so the total angular momentum $\lVert\mathbf{h}_{\text{tot}}\rVert$ is conserved, but energy is dissipated whenever the body nutates ($\boldsymbol{\omega}_d \neq \boldsymbol{\omega}$). For fixed $\lVert\mathbf{h}_{\text{tot}}\rVert$, the energy of pure spin about principal axis $i$ is
\[
T_i = \frac{\lVert\mathbf{h}_{\text{tot}}\rVert^2}{2\left(J_i + J_d\right)},
\]
so minor-axis spin is the \emph{maximum}-energy state and major-axis spin is the \emph{minimum}. Any small nutation about the minor axis loses energy, which carries the body further from minor-axis spin, so the nutation grows until the body reaches the only state the damper can't drain: major-axis spin.
\end{answer}

\item What happens to the rate of energy dissipation in the limit $C_d \to 0$? What about the limit $C_d \to \infty$?

\begin{answer}
\textbf{It goes to zero in both limits.} As $C_d \to 0$, the damper decouples from the body. The relative rate $\boldsymbol{\omega}_d - \boldsymbol{\omega}$ may be large, but $\dot T \to 0$ because of the factor of $C_d$. As $C_d \to \infty$, the damper locks to the body, so the relative rate shrinks like $1/C_d$ and
\[
\dot T \sim C_d \cdot \left(\frac{1}{C_d}\right)^2 = \frac{1}{C_d} \to 0 .
\]
The system then behaves like a single rigid body with inertia $\mathbf{J} + J_d\mathbf{I}$, which conserves energy. So an intermediate value of $C_d$ dissipates energy fastest, and a Kane damper must be tuned.
\end{answer}

\item A body with principal inertia matrix $\mathbf{J} = \operatorname{diag}(2,\ 3,\ 4)$ contains a Kane damper with $J_d = 1$. It starts in an equilibrium spin at rate $\Omega$ about $\hat b_1$ (the minor axis), with the damper co-rotating, and is then given a tiny disturbance. About which axis, and at what angular velocity, does the body eventually spin?

\begin{answer}
\textbf{About the major axis $\hat b_3$, with ${}^{B}\boldsymbol{\omega} = \pm\tfrac{3}{5}\,\Omega\,\hat b_3$.} Angular momentum, not energy, is conserved. Initially, the body and the damper both spin at $\Omega$ about $\hat b_1$, so
\[
\lVert\mathbf{h}_{\text{tot}}\rVert = \left(J_1 + J_d\right)\Omega = (2 + 1)\,\Omega = 3\,\Omega .
\]
In the final state, the body and the damper both spin at $\omega_3$ about $\hat b_3$, so
\[
\left(J_3 + J_d\right)\lvert\omega_3\rvert = 3\,\Omega
\;\Longrightarrow\;
\lvert\omega_3\rvert = \frac{3\,\Omega}{4 + 1} = \frac{3}{5}\,\Omega .
\]
The sign depends on the details of the disturbance, so either sign is correct. Physically, the body turns over so that $\hat b_3$ lines up with the inertially fixed angular momentum vector, along the original direction of $\hat b_1$. The kinetic energy drops by a factor of $(J_1 + J_d)/(J_3 + J_d) = 3/5$.
\end{answer}

\end{enumerate}

\end{document}
